% GENERATED FILE. Edit canonical lessons, then run npm run generate:books.
% canonical-source-hash: fnv1a64:f2a6f1dcaeb5965e
% canonical-lessons: MW-C282-sut, MW-C282-charkho, MW-C282-kaatno, MW-C282-mala, MW-C282-guthno

\chapter{Spinning and Beads}
\label{ch:mw-spinning-and-beads}
\hlchaptermodality{282}

\begin{chapteropening}
\textbf{By the end of this chapter:} \emph{I can say how thread is spun and beads are strung.}

Complete the last lesson of chapter 282: I can say how thread is spun and beads are strung.
\end{chapteropening}

\section[\texorpdfstring{\mw{सूत} (sūt)}{सूत (sūt)}]{\mw{सूत} (sūt) — yarn}

\label{lesson:MW-C282-sut}



\begin{quote}
Before the new one: say the Marwadi for to carve, then the Marwadi for a kiln.
\end{quote}

\subsection*{You'll want to know: \mw{सूत}}
\textbf{\mw{सूत}} — \emph{sūt} — \textquotedblleft{}yarn\textquotedblright{}.

\begin{quote}
\mw{यो} \mw{सूत} \mw{लाल} \mw{है।}
\end{quote}

\begin{quote}
\emph{yo sūt lāl hai.}
\end{quote}

\begin{quote}
This yarn is red.
\end{quote}

A masculine noun.

\begin{cousinweb}
\textbf{\mw{सूत}} (\emph{sūt}) is inherited from Sanskrit \emph{sūtra}, \textquotedblleft{}thread, cord\textquotedblright{}; CDIAL gives no Marwari form in this entry, but lists Hindi \emph{sūt}, \textquotedblleft{}thread\textquotedblright{} (CDIAL 13561).
\end{cousinweb}

\begin{grammarlens}[title={What you've built}]
One more word for this chapter.
\end{grammarlens}

\subsection*{Your turn}
\begin{itemize}
\raggedright
  \item \emph{Say it:} \emph{sūt}
  \item \emph{Say it:} \emph{sūt}, once more
  \item \emph{Say it:} \emph{yo sūt lāl hai.} — This yarn is red.
  \item \emph{Recall:} say the Marwadi for to carve, then the Marwadi for a kiln, then say \emph{sūt} again
\end{itemize}

\begin{tcolorbox}[breakable,colback=teal!4,colframe=teal!35!black,title={Before you move on}]
What does \mw{सूत} mean? (\textquotedblleft{}Yarn\textquotedblright{}.) Say it once more.
\end{tcolorbox}
\section[\texorpdfstring{\mw{चरखो} (charkho)}{चरखो (charkho)}]{\mw{चरखो} (charkho) — a spinning wheel}

\label{lesson:MW-C282-charkho}



\begin{quote}
Before the new one: say the Marwadi for a kiln, then the Marwadi for yarn.
\end{quote}

\subsection*{You'll want to know: \mw{चरखो}}
\textbf{\mw{चरखो}} — \emph{charkho} — \textquotedblleft{}a spinning wheel\textquotedblright{}.

\begin{quote}
\mw{दादी} \mw{कनै} \mw{एक} \mw{चरखो} \mw{है।}
\end{quote}

\begin{quote}
\emph{dādī kanai ek charkho hai.}
\end{quote}

\begin{quote}
Grandmother has a spinning wheel.
\end{quote}

A masculine noun.

\begin{cousinweb}
\textbf{\mw{चरखो}} (\emph{charkho}) is usually traced to Persian \emph{charkh}, \textquotedblleft{}a wheel\textquotedblright{}; Marwadi gives it the -\mw{ो} ending of a masculine noun.
\end{cousinweb}

\begin{grammarlens}[title={What you've built}]
One more word for this chapter.
\end{grammarlens}

\subsection*{Your turn}
\begin{itemize}
\raggedright
  \item \emph{Say it:} \emph{charkho}
  \item \emph{Say it:} \emph{charkho}, once more
  \item \emph{Say it:} \emph{dādī kanai ek charkho hai.} — Grandmother has a spinning wheel.
  \item \emph{Recall:} say the Marwadi for a kiln, then the Marwadi for yarn, then say \emph{charkho} again
\end{itemize}

\begin{tcolorbox}[breakable,colback=teal!4,colframe=teal!35!black,title={Before you move on}]
What does \mw{चरखो} mean? (\textquotedblleft{}A spinning wheel\textquotedblright{}.) Say it once more.
\end{tcolorbox}
\section[\texorpdfstring{\mw{कातणो} (kātṇo)}{कातणो (kātṇo)}]{\mw{कातणो} (kātṇo) — to spin thread}

\label{lesson:MW-C282-kaatno}



\begin{quote}
Before the new one: say the Marwadi for yarn, then the Marwadi for a spinning wheel.
\end{quote}

\subsection*{You'll want to know: \mw{कातणो}}
\textbf{\mw{कातणो}} — \emph{kātṇo} — \textquotedblleft{}to spin thread\textquotedblright{}.

\begin{quote}
\mw{दादी} \mw{रोज} \mw{सूत} \mw{कातै} \mw{है।}
\end{quote}

\begin{quote}
\emph{dādī roj sūt kātai hai.}
\end{quote}

\begin{quote}
Grandmother spins yarn every day.
\end{quote}

Past (I …): \mw{म्हैं} \mw{सूत} \mw{कात्यो}, \emph{mhaĩ sūt kātyo}.

\begin{cousinweb}
\textbf{\mw{कातणो}} (\emph{kātṇo}) is inherited from a reconstructed \emph{*kartati}, \textquotedblleft{}spins\textquotedblright{}; CDIAL lists Old Marwari \emph{kātaï}, \textquotedblleft{}spins\textquotedblright{} (CDIAL 2855).
\end{cousinweb}

\begin{grammarlens}[title={What you've built}]
One more word for this chapter.
\end{grammarlens}

\subsection*{Your turn}
\begin{itemize}
\raggedright
  \item \emph{Say it:} \emph{kātṇo}
  \item \emph{Say it:} \emph{kātṇo}, once more
  \item \emph{Say it:} \emph{dādī roj sūt kātai hai.} — Grandmother spins yarn every day.
  \item \emph{Recall:} say the Marwadi for yarn, then the Marwadi for a spinning wheel, then say \emph{kātṇo} again
\end{itemize}

\begin{tcolorbox}[breakable,colback=teal!4,colframe=teal!35!black,title={Before you move on}]
What does \mw{कातणो} mean? (\textquotedblleft{}To spin thread\textquotedblright{}.) Say it once more.
\end{tcolorbox}
\section[\texorpdfstring{\mw{माळा} (māḷā)}{माळा (māḷā)}]{\mw{माळा} (māḷā) — a garland; a string of beads}

\label{lesson:MW-C282-mala}



\begin{quote}
Before the new one: say the Marwadi for a spinning wheel, then the Marwadi for to spin thread.
\end{quote}

\subsection*{You'll want to know: \mw{माळा}}
\textbf{\mw{माळा}} — \emph{māḷā} — \textquotedblleft{}a garland; a string of beads\textquotedblright{}.

\begin{quote}
\mw{म्हारी} \mw{माळा} \mw{नुवी} \mw{है।}
\end{quote}

\begin{quote}
\emph{mhārī māḷā nuvī hai.}
\end{quote}

\begin{quote}
My garland is new.
\end{quote}

A feminine noun.

\begin{cousinweb}
\textbf{\mw{माळा}} (\emph{māḷā}) is inherited from Sanskrit \emph{mālā}, \textquotedblleft{}a wreath, a garland\textquotedblright{}, which CDIAL takes from Dravidian; it lists Old Marwari \emph{māla}, \textquotedblleft{}garland\textquotedblright{} (CDIAL 10092).
\end{cousinweb}

\begin{grammarlens}[title={What you've built}]
One more word for this chapter.
\end{grammarlens}

\subsection*{Your turn}
\begin{itemize}
\raggedright
  \item \emph{Say it:} \emph{māḷā}
  \item \emph{Say it:} \emph{māḷā}, once more
  \item \emph{Say it:} \emph{mhārī māḷā nuvī hai.} — My garland is new.
  \item \emph{Recall:} say the Marwadi for a spinning wheel, then the Marwadi for to spin thread, then say \emph{māḷā} again
\end{itemize}

\begin{tcolorbox}[breakable,colback=teal!4,colframe=teal!35!black,title={Before you move on}]
What does \mw{माळा} mean? (\textquotedblleft{}A garland; a string of beads\textquotedblright{}.) Say it once more.
\end{tcolorbox}
\section[\texorpdfstring{\mw{गूंथणो} (gū̃thṇo)}{गूंथणो (gū̃thṇo)}]{\mw{गूंथणो} (gū̃thṇo) — to string beads; to plait}

\label{lesson:MW-C282-guthno}



\begin{quote}
Before the new one: say the Marwadi for to spin thread, then the Marwadi for a garland.
\end{quote}

\subsection*{You'll want to know: \mw{गूंथणो}}
\textbf{\mw{गूंथणो}} — \emph{gū̃thṇo} — \textquotedblleft{}to string beads; to plait\textquotedblright{}.

\begin{quote}
\mw{छोरी} \mw{माळा} \mw{गूंथै} \mw{है।}
\end{quote}

\begin{quote}
\emph{chhorī māḷā gū̃thai hai.}
\end{quote}

\begin{quote}
The girl strings a garland.
\end{quote}

Past (I …): \mw{म्हैं} \mw{माळा} \mw{गूंथी}, \emph{mhaĩ māḷā gū̃thī}.

\begin{cousinweb}
\textbf{\mw{गूंथणो}} (\emph{gū̃thṇo}) is inherited from Sanskrit \emph{guphati}, \textquotedblleft{}strings together\textquotedblright{}; CDIAL lists Old Marwari \emph{guthaï}, \textquotedblleft{}strings, threads\textquotedblright{} (CDIAL 4205).
\end{cousinweb}

\begin{grammarlens}[title={What you've built}]
One more word for this chapter.
\end{grammarlens}

\subsection*{Your turn}
\begin{itemize}
\raggedright
  \item \emph{Say it:} \emph{gū̃thṇo}
  \item \emph{Say it:} \emph{gū̃thṇo}, once more
  \item \emph{Say it:} \emph{chhorī māḷā gū̃thai hai.} — The girl strings a garland.
  \item \emph{Recall:} say the Marwadi for to spin thread, then the Marwadi for a garland, then say \emph{gū̃thṇo} again
\end{itemize}

\begin{tcolorbox}[breakable,colback=teal!4,colframe=teal!35!black,title={Before you move on}]
What does \mw{गूंथणो} mean? (\textquotedblleft{}To string beads; to plait\textquotedblright{}.) Say it once more.
\end{tcolorbox}
